\documentclass[11pt,a4paper]{amsart}

\usepackage[margin=1in]{geometry}
\usepackage{amsmath,amssymb,amsthm}
\usepackage[hidelinks]{hyperref}
\emergencystretch=1em

\hypersetup{
  pdftitle={A Cauchy Limit for a Sawtooth Sum},
  pdfauthor={Sangyoon Kwon},
  pdfsubject={A solution to Erdos Problem 1002},
  pdfkeywords={continued fractions, discrepancy sums, Cauchy law,
    stable limits, Gauss map, Erdos Problem 1002}
}

\newtheorem{theorem}{Theorem}[section]
\newtheorem{proposition}[theorem]{Proposition}
\newtheorem{lemma}[theorem]{Lemma}
\newtheorem{corollary}[theorem]{Corollary}
\theoremstyle{definition}
\newtheorem{definition}[theorem]{Definition}
\theoremstyle{remark}
\newtheorem{remark}[theorem]{Remark}

\newcommand{\T}{\mathbb T}
\newcommand{\N}{\mathbb N}
\newcommand{\Z}{\mathbb Z}
\newcommand{\R}{\mathbb R}
\newcommand{\E}{\mathbb E}
\newcommand{\Pp}{\mathbb P}
\newcommand{\ind}{\mathbf 1}
\newcommand{\dist}{\operatorname{dist}}
\newcommand{\Var}{\operatorname{Var}}
\newcommand{\Cov}{\operatorname{Cov}}
\newcommand{\Leb}{\operatorname{Leb}}
\newcommand{\e}{\mathrm e}

\title{A Cauchy Limit for a Sawtooth Sum}
\author{Sangyoon Kwon}
\address{Independent researcher}
\email{is011061@gmail.com}
\date{Version 8, August 6, 2026}
\subjclass[2020]{11K38, 11A55, 37A50, 60F05}
\keywords{continued fractions, discrepancy sums, Cauchy law, stable limits,
Gauss map}

\begin{document}
\begin{abstract}
For $S_n(\alpha)=\sum_{k\le n}(1/2-\{k\alpha\})$, with $\alpha$ uniformly distributed on $(0,1)$, we prove that $S_n(\alpha)/\log n$ converges in distribution to the centered Cauchy law of scale $1/(2\pi)$.  An exact Euclidean-algorithm reciprocity formula reduces the sum to an alternating continued-fraction cost.  Its principal part is a Gauss digit marked by a rapidly oscillating torus coordinate, while the carry remainder has uniformly bounded summands.  Local cylinder oscillation and a resonance decomposition yield a Poisson limit for the large jumps and a uniform variance bound for the compensated small jumps; a pullback carry construction gives a weak law for the bounded remainder.  Consequently the limiting distribution function is $\frac12+\pi^{-1}\arctan(2\pi c)$.
\end{abstract}

\maketitle

\section{Introduction and statement of the result}

Write $\{x\}=x-\lfloor x\rfloor$ and let logarithms be natural.  For
$n\in\N_0$, put
\[
 S_n(\alpha)=\sum_{k=1}^{n}\left(\frac12-\{k\alpha\}\right),
 \qquad 0<\alpha<1,
\]
with the empty-sum convention $S_0=0$.  For $n\ge1$, the question is whether
$S_n(\alpha)/\log n$, regarded as a random variable on
the Lebesgue probability space $(0,1)$, has a limiting law.  This is
Erd\H{o}s Problem~\#1002, originating in \cite{Er64b}; see also
\cite{Bloom1002}.  Erd\H{o}s used the opposite sign convention; since
$\alpha\mapsto1-\alpha$ preserves Lebesgue measure and reverses the
sum for irrational $\alpha$, the two distributional formulations are
equivalent.  Kesten's two-parameter result \cite{Ke60} provides the
closely related averaged-origin case; Borda \cite{Borda25} computed the
normalizing constant in that law for the sawtooth observable.  This is an
annealed spatial law, averaging both the rotation number and the starting
point.  Dolgopyat and Sarig \cite{DS20} proved an annealed temporal Cauchy
law in which the rotation number and the time are randomized.  Neither
ensemble is the one considered here: the time $n$ and the starting point are
fixed, and only $\alpha$ is averaged.  A related Cauchy limit phenomenon
occurs in Vardi's theorem for Dedekind sums \cite{Vardi93}, although the
observable and averaging ensemble are different from those considered here.
The answer to this
fixed-origin, fixed-time problem is as follows.

\begin{theorem}\label{thm:main}
As $n\to\infty$,
\[
 \frac{S_n(\alpha)}{\log n}
 \ \Longrightarrow\
 \operatorname{Cauchy}\!\left(0,\frac1{2\pi}\right).
\]
Equivalently, for every $c\in\R$,
\[
 \lim_{n\to\infty}\Leb\left\{\alpha\in(0,1):
 \frac{S_n(\alpha)}{\log n}\le c\right\}
 =\frac12+\frac1\pi\arctan(2\pi c).
\]
\end{theorem}

Here $\operatorname{Cauchy}(0,\gamma)$ denotes the centered Cauchy law with
characteristic function $t\mapsto\exp(-\gamma|t|)$.

The proof has four parts.  First an exact reciprocity identity expresses
$S_n$ as an alternating sum along the Gauss map.  Second, continued-fraction
coordinates separate each summand into a heavy-tailed marked digit and a
bounded carry remainder.  Third, the marked digits satisfy a Poisson limit;
the main point is that the initial measure lies on the oscillating graph
$\alpha\mapsto(\alpha,\{n\alpha\})$, so ordinary Gauss mixing is insufficient.
We instead integrate Fourier modes on complete continued-fraction cylinders,
before and after their unique resonance scale.  Finally, the carry is realized
as a measurable graph over the natural extension of a Gauss--torus cocycle;
a reset argument and a two-block correlation estimate give a weak law for the
bounded remainder.

Throughout, constants denoted by $C,c$ may change from line to line.  A bound
$O_A(L^{-A})$ is uniform in the time indices satisfying the displayed
hypotheses.  We put $L=\log n$ whenever $n$ is fixed.

\section{Exact renormalization and continued-fraction coordinates}

Let $T:(0,1)\to(0,1)$ be the Gauss map,
\[
 Tx=\left\{\frac1x\right\},
 \qquad a_1(x)=\left\lfloor\frac1x\right\rfloor.
\]
Rational points will always be ignored, since they form a null set.

\begin{proposition}[Reciprocity]\label{prop:reciprocity}
For irrational $x\in(0,1)$ and an integer $N\ge1$, set
\[
 m=\lfloor Nx\rfloor,\qquad u=\{Nx\}.
\]
Then
\begin{equation}\label{eq:reciprocity}
 S_N(x)=-S_m(Tx)+\Phi(x,u),
 \qquad
 \Phi(x,u)=\frac{u(1-u)}{2x}-\frac u2.
\end{equation}
\end{proposition}

\begin{proof}
Since $x$ is irrational,
\[
 S_N(x)=\frac N2-\frac{xN(N+1)}2+\sum_{k=1}^{N}\lfloor kx\rfloor.
\]
Counting lattice points below $j=kx$ in the other order gives
\[
 \sum_{k=1}^{N}\lfloor kx\rfloor
 =Nm-\sum_{j=1}^{m}\left\lfloor\frac jx\right\rfloor.
\]
If $a=\lfloor1/x\rfloor$, then $1/x=a+Tx$, and hence
\[
 \sum_{j=1}^{m}\left\lfloor\frac jx\right\rfloor
 =\frac{am(m+1)}2+\sum_{j=1}^{m}\lfloor jTx\rfloor.
\]
Substitution and the identity $m=Nx-u$ reduce the remaining polynomial
expression to $u(1-u)/(2x)-u/2$.
\end{proof}



Starting with $N_0=n$, $x_0=\alpha$, define
\begin{equation}\label{eq:descent}
 N_{j+1}=\lfloor N_jx_j\rfloor,
 \qquad x_{j+1}=Tx_j,
 \qquad u_j=\{N_jx_j\}.
\end{equation}
Let $\tau_n=\min\{j:N_j=0\}$.  Iterating
Proposition~\ref{prop:reciprocity} gives the exact identity
\begin{equation}\label{eq:exact-sum}
 S_n(\alpha)=\sum_{j=0}^{\tau_n-1}(-1)^j\Phi(x_j,u_j).
\end{equation}

Write
\[
 \alpha=[0;a_1,a_2,\ldots],\qquad x_j=T^j\alpha,
 \qquad a_{j+1}=a_1(x_j),
\]
and let $p_j/q_j$ be the convergents, normalized by
\[
 p_{-1}=1,\quad p_0=0,\qquad q_{-1}=0,\quad q_0=1.
\]
The standard identities
\begin{equation}\label{eq:beta}
 \beta_j:=x_0x_1\cdots x_j
 =(-1)^j(q_j\alpha-p_j)
 =\frac1{q_{j+1}+q_jx_{j+1}}
\end{equation}
will be used repeatedly; set $\beta_{-1}=1$.

Define the continuous heights and their errors by
\[
 C_j=n\beta_{j-1},\qquad E_j=C_j-N_j.
\]
Then $C_{j+1}=x_jC_j$, and \eqref{eq:descent} yields
\begin{equation}\label{eq:error-recurrence}
 E_{j+1}=x_jE_j+u_j,\qquad E_0=0.
\end{equation}
For every $r,\ell\ge0$,
\begin{equation}\label{eq:fibonacci-product}
 x_rx_{r+1}\cdots x_{r+\ell-1}\le F_{\ell+1}^{-1},
\end{equation}
where $F_1=F_2=1$.  Indeed, the reciprocal of the product is a
continuant with all digits at least one.  Iterating \eqref{eq:error-recurrence}
therefore gives
\begin{equation}\label{eq:E-bound}
 0\le E_j\le E_*:=\sum_{\ell\ge1}F_\ell^{-1}<\infty
\end{equation}
uniformly in $j,n,\alpha$.

Put
\begin{equation}\label{eq:theta}
 \theta_j=\{n\beta_j\}=\{(-1)^jnq_j\alpha\},
 \qquad \theta_{-1}=0.
\end{equation}
From $q_{j+1}=a_{j+1}q_j+q_{j-1}$ we obtain
\begin{equation}\label{eq:theta-cocycle}
 \theta_{j+1}=\{\theta_{j-1}-a_{j+1}\theta_j\}.
\end{equation}
Since $\{E_j\}=\theta_{j-1}$, write
\[
 E_j=d_j+\theta_{j-1},\qquad d_j=\lfloor E_j\rfloor.
\]
The equality $C_{j+1}=N_{j+1}+u_j+x_jE_j$ implies
\begin{equation}\label{eq:u-carry}
 u_j=\{\theta_j-x_j(d_j+\theta_{j-1})\}.
\end{equation}

Let
\[
 W(t)=\frac{\{t\}(1-\{t\})}{2},\qquad t\in\R/\Z.
\]
This function is $1/2$-Lipschitz and $0\le W\le1/8$.

\begin{proposition}[Principal term]\label{prop:principal}
There is an absolute constant $C_0$ such that, for every $j$,
\begin{equation}\label{eq:principal-split}
 \Phi(x_j,u_j)=a_{j+1}W(\theta_j)+B_j,
 \qquad |B_j|\le C_0.
\end{equation}
\end{proposition}

\begin{proof}
By \eqref{eq:u-carry} and \eqref{eq:E-bound},
\[
 \dist_{\T}(u_j,\theta_j)\le x_jE_*.
\]
Since $x_j^{-1}=a_{j+1}+x_{j+1}$,
\[
 \Phi(x_j,u_j)-a_{j+1}W(\theta_j)
 =a_{j+1}\bigl(W(u_j)-W(\theta_j)\bigr)
  +x_{j+1}W(u_j)-\frac{u_j}{2}.
\]
The first term has absolute value at most
$a_{j+1}x_jE_*/2\le E_*/2$, and the other two are bounded.
\end{proof}


\section{Gauss estimates and local oscillation}

The Gauss invariant probability measure is
\begin{equation}\label{eq:gauss-measure}
 d\nu(x)=h(x)\,dx,
 \qquad h(x)=\frac1{\log2\,(1+x)}.
\end{equation}
The constant
\begin{equation}\label{eq:lambda}
 \lambda=\int_0^1-\log x\,d\nu(x)
 =\frac{\pi^2}{12\log2}
\end{equation}
is the L\'evy constant.  Since $\log|T'(x)|=-2\log x$, the Lyapunov
exponent of the Gauss map is $2\lambda$.

For a word $w=(a_1,\ldots,a_r)$ let $I_w$ be the corresponding
continued-fraction cylinder.  If $p_r/q_r$ is its last convergent, the inverse
branch is
\[
 h_w(y)=\frac{p_r+p_{r-1}y}{q_r+q_{r-1}y},
 \qquad
 |h_w'(y)|=(q_r+q_{r-1}y)^{-2},
\]
and
\begin{equation}\label{eq:cylinder-length}
 |I_w|=\frac1{q_r(q_r+q_{r-1})}.
\end{equation}
Conditionally on $I_w$, the density of $T^r\alpha=y$ is
\[
 \rho_w(y)=\frac{q_r(q_r+q_{r-1})}{(q_r+q_{r-1}y)^2}.
\]
Thus $\sup_w\|\rho_w\|_{BV}<\infty$.

We use the following standard consequences of the Lasota--Yorke inequality
for the Gauss transfer operator; see, for example, \cite{Baladi00}.  They are
recorded with the uniformity needed below.

\begin{lemma}[Gauss estimates]\label{lem:gauss-estimates}
There are $C,c>0$ and $0<\rho<1$ such that:
\begin{enumerate}
\item for every integer $r\ge0$ and $f\in BV(0,1)$,
\[
 \|\mathcal L^rf-h\!\int f\|_1\le C\rho^r\|f\|_{BV};
\]
\item for an integer $A\ge1$, conditionally on any past cylinder,
\[
 \Pp(a_{j+1}\ge A\mid a_1,\ldots,a_j)\le C/A;
\]
in particular, for distinct $j_1<\cdots<j_s$ and positive integers
$A_1,\ldots,A_s$,
\begin{equation}\label{eq:multi-digit-tail}
 \Pp(a_{j_i+1}\ge A_i,\ 1\le i\le s)
 \le C^s\prod_{i=1}^s A_i^{-1};
\end{equation}
\item uniformly for $r\ge1$ and $v>0$,
\begin{equation}\label{eq:q-large-deviation}
 \Pp(|\log q_r-\lambda r|>v)
 \le C\exp\{-c\min(v^2/r,v)\}.
\end{equation}
\end{enumerate}
\end{lemma}

\begin{proof}
For completeness, the transfer operator is
\[
 \mathcal Lf(x)=\sum_{a\ge1}(a+x)^{-2}
 f\!\left((a+x)^{-1}\right).
\]
On two-step inverse branches the contraction is at most $1/2$; the sums of
the inverse Jacobians and of their variations converge.  The branchwise chain
rule gives
\[
 \Var(\mathcal L^2f)\le\rho_0\Var(f)+C\|f\|_1,
 \qquad \rho_0<1.
\]
Helly compactness and exactness yield (i).  The conditional densities
$\rho_w$ have uniformly bounded variation, and
$\{a_1\ge A\}=(0,1/A]$ up to endpoints; this proves (ii), and iteration proves
\eqref{eq:multi-digit-tail}.  Finally, the perturbed operators
\[
 \mathcal L_t f(x)=\sum_{a\ge1}(a+x)^{-2+t}
 f\!\left((a+x)^{-1}\right)
\]
form an analytic family on $BV$ for $|t|<1/2$.  Indeed, on every smaller
complex disc the branch series and all its $t$-derivatives converge in
operator norm, since their bounds reduce to convergent series of the form
$\sum_{a\ge1}a^{-2+t_0}(1+(\log a)^k)$.  Analytic perturbation of the simple
leading eigenvalue \cite[Chap.~VII]{Kato95} therefore gives, after decreasing
$t_0>0$ if necessary,
\[
 \begin{aligned}
  \mathcal L_t&=\kappa(t)\Pi_t+N_t,
  &\Pi_tN_t&=N_t\Pi_t=0,\\
  \sup_{|t|\le t_0}\|\Pi_t\|_{BV\to BV}&<\infty,
  &\sup_{|t|\le t_0}\|N_t^r\|_{BV\to L^1}&\le C\rho_1^r
  \quad(r\ge0).
 \end{aligned}
\]
where these bounds are for real $t$ and
$0<\rho_1<\min_{|t|\le t_0}\kappa(t)$.  Moreover,
\[
 \kappa'(0)=\int_0^1-\log x\,d\nu(x)=\lambda,
 \qquad
 \E\exp\!\left(t\sum_{i<r}-\log T^i\alpha\right)
 =\int_0^1\mathcal L_t^r\ind\,dx.
\]
Consequently,
\[
 \log\kappa(t)=\lambda t+O(t^2),
 \qquad
 \E\exp\!\left(t\sum_{i<r}-\log T^i\alpha\right)
 \le C\kappa(t)^r
\]
for $|t|$ sufficiently small; the same estimate holds for negative $t$.
For $0<v\le r$, Chernoff's inequality with $|t|\asymp v/r$ gives
$C\exp(-cv^2/r)$, while for $v>r$ a fixed sufficiently small $|t|$ gives
$C\exp(-cv)$.  Thus the asserted Bernstein bound holds for
$\sum_{i<r}-\log T^i\alpha$.
Since
\[
 \left|\log q_r-\sum_{i<r}-\log T^i\alpha\right|\le\log2,
\]
we obtain (iii).
\end{proof}

\begin{lemma}[Conditional multi-block mixing]\label{lem:conditional-mixing}
Let $s,M\ge1$ and $d\ge0$ be integers, let $I_w$ be a cylinder of depth $d$,
and let $t_1,\ldots,t_s$ be integers satisfying
$d+M\le t_1<\cdots<t_s$ and $t_{i+1}-t_i\ge M$.  If
$g_1,\ldots,g_s\in BV(0,1)$, then
\begin{equation}\label{eq:conditional-mixing}
 \left|
 \E\left(\prod_{i=1}^s g_i(T^{t_i}\alpha)\,\middle|\,I_w\right)
 -\prod_{i=1}^s\int g_i\,d\nu
 \right|
 \le C_s\rho^M\prod_{i=1}^s\|g_i\|_{BV}.
\end{equation}
The constants are uniform in the cylinder $I_w$.
\end{lemma}

\begin{proof}
Conditioned on $I_w$, the density of $T^d\alpha$ is $\rho_w$, whose $BV$
norm is uniformly bounded by the displayed formula for $\rho_w$ above.  More
precisely, change of variables by the inverse branch of $I_w$ gives the exact
cylinder-transfer identity
\[
 \E\left(\prod_{i=1}^s g_i(T^{t_i}\alpha)\,\middle|\,I_w\right)
 =\int_0^1\prod_{i=1}^s g_i(T^{t_i-d}y)\rho_w(y)\,dy.
\]
Apply Lemma~\ref{lem:gauss-estimates}(i) across the first gap, multiply by the
next $BV$ observable, and iterate.  The product inequality
$\|fg\|_{BV}\le C\|f\|_{BV}\|g\|_{BV}$ and the fixed number $s$ of factors
give \eqref{eq:conditional-mixing}.
\end{proof}

Fix henceforth
\begin{equation}\label{eq:scales}
 L=\log n,\qquad H=L^{3/4},\qquad
 m_n=\left\lfloor L/\lambda\right\rfloor,
\end{equation}
and define the bulk set
\begin{equation}\label{eq:bulk}
 \mathcal J_n=\{j\in\Z:200H\le j\le m_n-200H\}.
\end{equation}
By \eqref{eq:q-large-deviation}, for every fixed $\delta>0$ and
every integer $0\le t\le2m_n$,
\begin{equation}\label{eq:local-q-good}
 \Pp(q_t\notin[\e^{\lambda t-\delta H},
                    \e^{\lambda t+\delta H}])
 \le C_\delta \e^{-c_\delta L^{1/2}}.
\end{equation}
We always use \eqref{eq:local-q-good} as a union of complete depth-$t$
cylinders.  We never insert a global good-event indicator inside an
oscillatory integral.

The first ingredient controls local cancellation of the two continuants in a
torus character.

\begin{lemma}[Shrinking anti-concentration]\label{lem:anti}
For every $(r,s)\in\Z^2\setminus\{(0,0)\}$, every $j\ge1$, and every
$\eta>0$,
\begin{equation}\label{eq:anti}
 \Pp(|sq_j-rq_{j-1}|<\eta q_j)
 \le 4\eta+4F_{j+1}^{-2}
 \le C\bigl(\eta+\e^{-2(\log\varphi)j}\bigr),
 \qquad \varphi=(1+\sqrt5)/2.
\end{equation}
\end{lemma}

\begin{proof}
Put $Y_j=q_{j-1}/q_j=[0;a_j,\ldots,a_1]$.  For a word
$w=(a_1,\ldots,a_j)$, let $w^{\rm rev}=(a_j,\ldots,a_1)$.  Symmetry of the
continuant gives
\[
 q_j(w^{\rm rev})=q_j(w),
 \qquad p_j(w^{\rm rev})=q_{j-1}(w),
 \qquad q_{j-1}(w^{\rm rev})=p_j(w).
\]
Consequently $Y_j(w)=p_j(w^{\rm rev})/q_j(w^{\rm rev})$, while
\eqref{eq:cylinder-length} gives the uniform distortion comparison
\[
 \frac{|I_w|}{|I_{w^{\rm rev}}|}
 =\frac{q_j(w)+p_j(w)}{q_j(w)+q_{j-1}(w)}\le2.
\]
The rational $p_j(w^{\rm rev})/q_j(w^{\rm rev})$ is an endpoint of
$I_{w^{\rm rev}}$, and $q_j\ge F_{j+1}$.  Thus, if $I\subset(0,1)$ is an
interval and $\delta_j=F_{j+1}^{-2}$, every reversed cylinder contributing to
$\{Y_j\in I\}$ lies in the $\delta_j$-neighborhood $I^{(\delta_j)}$.  Since
word reversal is a bijection and the depth-$j$ cylinders are disjoint up to
endpoints,
\[
 \Pp(Y_j\in I)
 \le2\sum_{w:Y_j(w)\in I}|I_{w^{\rm rev}}|
 \le2|I^{(\delta_j)}|
 \le2|I|+4\delta_j.
\]
If $r\ne0$, the event in \eqref{eq:anti} is $|s-rY_j|<\eta$ and its possible
$Y_j$-values lie in an interval of length at most $2\eta/|r|\le2\eta$.
If $r=0$, then $s\ne0$; the event is empty when $\eta\le|s|$, while otherwise
its probability is $1\le\eta$.  This proves the first bound uniformly in
$(r,s)$.  The second follows from $F_{j+1}^{-2}\ll\varphi^{-2j}$.
\end{proof}

The second ingredient is a complete-cylinder oscillatory estimate.  Its
formulation makes the measurability needed for integration by parts explicit.

\begin{lemma}[Descendant-cylinder estimate]\label{lem:descendant}
Let $\varepsilon>0$, and let $\{I_w\}$ be cylinders of a fixed depth $d$.
On $I_w$ let
$Q_w\in\Z\setminus\{0\}$ be fixed.  Let $k>d$ and let $A$ be constant on every
depth-$k$ cylinder, with $|A|\le1$.  Suppose that, below $I_w$, the support of
$A$ is a union of depth-$k$ descendants satisfying $q_k\le R_w$, where
$R_w>0$ and
\[
 R_w^2\le\varepsilon n|Q_w|.
\]
Then
\begin{equation}\label{eq:descendant}
 \left|\sum_w\int_{I_w}
       \e^{2\pi i nQ_w\alpha}A(\alpha)\,d\alpha\right|
 \le C\varepsilon.
\end{equation}
Conversely, let $Q\in\Z\setminus\{0\}$.  If a depth-$k$ cylinder $J$
satisfies $q_k^2\ge\varepsilon^{-1}n|Q|$, then
\begin{equation}\label{eq:phase-freeze}
 \sup_{\alpha,\alpha'\in J}|nQ(\alpha-\alpha')|
 \le C\varepsilon.
\end{equation}
\end{lemma}

\begin{proof}
Let $q(w)$ and $q_-(w)$ be the last two denominators of the prefix $w$.
If $(\ell,\ell')$ is the denominator pair of a continuation, the final
denominator is
$q_k=q(w)\ell+q_-(w)\ell'\ge q(w)\ell$.  The number of descendants with
$q_k\le R_w$ is therefore $O(R_w^2/q(w)^2)$.  On each complete descendant
cylinder $J$,
\[
 \left|\int_J\e^{2\pi i nQ_w\alpha}\,d\alpha\right|
 \le(\pi n|Q_w|)^{-1}.
\]
Thus the contribution below $w$ is at most $C\varepsilon/q(w)^2$.  By
\eqref{eq:cylinder-length}, $q(w)^{-2}\le2|I_w|$, and summation proves
\eqref{eq:descendant}.  The cylinder length bound $|J|\le q_k^{-2}$ gives
\eqref{eq:phase-freeze}.
\end{proof}


\section{Marked factorization and resonances}

For $D>0$ let $\mathcal P_D(L)$ consist of functions
\begin{equation}\label{eq:one-time-poly}
 F(a,\theta)=\sum_{|v|\le L^D}c(a,v)\e^{2\pi iv\theta},
 \qquad c(a,v)=0\quad(a>L^D),
\end{equation}
such that $\sum_{a,v}|c(a,v)|\le L^D$.  For times
$j_1<\cdots<j_r$ in $\mathcal J_n$, call the tuple good if
\begin{equation}\label{eq:good-tuple-separation}
 j_{\ell+1}-j_\ell\ge200H
\end{equation}
and, for every $i<\ell$,
\begin{equation}\label{eq:good-tuple-resonance}
 \left|j_\ell-\frac{m_n+j_i}{2}\right|\ge200H.
\end{equation}
The number of non-good $r$-tuples is $O_r(L^{r-1}H)$.

\begin{proposition}[Marked factorization]\label{prop:marked-factorization}
Fix an integer $r\ge1$ and $D,A>0$.  If $(j_1,\ldots,j_r)$ is good and
$F_1,\ldots,F_r\in\mathcal P_D(L)$, then
\begin{align}\label{eq:marked-factorization}
 &\int_0^1\prod_{\ell=1}^r
 F_\ell(a_{j_\ell+1},\theta_{j_\ell})\,d\alpha\\
 &\qquad=
 \prod_{\ell=1}^r\int_0^1\int_0^1
 F_\ell(a_1(x),\theta)\,d\theta\,d\nu(x)
 +O_{r,D,A}(L^{-A}).\notag
\end{align}
\end{proposition}

\begin{proof}
Expand both in the digit values and in the Fourier modes
$v_1,\ldots,v_r$, and pull out the scalar coefficient of each
digit--Fourier term.  The remaining amplitude is a product of cylinder
indicators and hence has modulus at most one; the total $\ell^1$-mass of the
extracted coefficients is at most $L^{Dr}$.  If every $v_\ell=0$, the
integrand is a product of digit functions.  For each such factor,
\[
 g_\ell(x)=\sum_a c_\ell(a,0)\ind_{\{a_1(x)=a\}},
 \qquad
 \Var(g_\ell)\le2\sum_a|c_\ell(a,0)|\le2L^D.
\]
This variation estimate comes from the coefficient $\ell^1$ bound in
\eqref{eq:one-time-poly}, and in particular $\|g_\ell\|_{BV}\ll L^D$; merely
counting its jumps would not give the stated bound.  Repeated use of
Lemma~\ref{lem:conditional-mixing} and
\eqref{eq:good-tuple-separation} factors its mean with error
$L^{O_{r,D}(1)}\rho^{200H}=O_A(L^{-A})$.

Suppose now that $v_s\ne0$ and $v_\ell=0$ for $\ell>s$.  By
\eqref{eq:theta}, the phase is
\[
 \e^{2\pi i nQ\alpha},
 \qquad
 Q=\sum_{\ell=1}^{s}v_\ell(-1)^{j_\ell}q_{j_\ell}.
\]
The continuants grow at least at the Fibonacci rate.  The separation
condition and $|v_\ell|\le L^D$ therefore give, deterministically,
\begin{equation}\label{eq:Q-dominance}
 \tfrac12q_{j_s}\le|Q|\le2L^Dq_{j_s}
\end{equation}
for all sufficiently large $n$.

The nominal resonance time is
\[
 R_s=\frac{m_n+j_s}{2}.
\]
Put $k_\pm=\lfloor R_s\pm40H\rfloor$.  We now make only local complete-cylinder
cuts.  At depth $j_s$ retain the cylinder union on which
\[
 \e^{\lambda j_s-\delta H}\le q_{j_s}
 \le \e^{\lambda j_s+\delta H}.
\]
At depths $k_-$ and $k_+$ retain, respectively,
\[
 q_{k_-}\le \e^{\lambda k_-+\delta H},\qquad
 q_{k_+}\ge \e^{\lambda k_+-\delta H},
\]
where $\delta>0$ is sufficiently small.  By
\eqref{eq:local-q-good}, the total discarded mass is
$O(\e^{-cL^{1/2}})$.  On every retained prefix and descendant,
\begin{equation}\label{eq:pre-post-resonance}
 q_{k_-}^2\le \e^{-cH}n|Q|,
 \qquad
 q_{k_+}^2\ge \e^{cH}n|Q|.
\end{equation}

The factors after $s$ are zero torus modes and hence digit functions.  By
\eqref{eq:good-tuple-resonance}, each lies either at least $100H$ before
$k_-$ or at least $100H$ after $k_+$.  For a fixed digit--Fourier term, write
$A_-$ for the product of all pre-resonance factors together with the retained
depth-$k_-$ cutoff, write $A_+$ for the product of the post-resonance digit
factors, denote those digit functions by $g_\ell$, and put
\[
 P_+=\prod_{\substack{\ell>s\\j_\ell\ge k_++100H}}
       \int g_\ell\,d\nu.
\]
Partition into complete depth-$k_+$ cylinders.  On every retained such
cylinder $J$, phase freezing and conditional mixing give
\[
 \int_J \e^{2\pi i nQ\alpha}A_-(\alpha)A_+(\alpha)\,d\alpha
 =P_+\int_J \e^{2\pi i nQ\alpha}A_-(\alpha)\,d\alpha
 +O\!\left(|J|L^{O_{r,D}(1)}
   (\e^{-cH}+\rho^{cH})\right).
\]
Here the phase-freezing error is \eqref{eq:phase-freeze}, and the mixing
error is summed by cylinder mass before the extracted coefficient
$\ell^1$-mass is restored.

Extend this stationary-mean replacement to the discarded depth-$k_+$
cylinders and restore them.  Their total mass is
$O(\e^{-cL^{1/2}})$, so this costs at most
$L^{O_{r,D}(1)}\e^{-cL^{1/2}}$ and removes every future cutoff.  The remaining
amplitude $A_-$ is constant on complete depth-$k_-$ descendants, its support
still satisfies the $q_{k_-}$ cutoff, and $Q$ is fixed on each
depth-$(j_s+1)$ prefix.  The first inequality in
\eqref{eq:pre-post-resonance} and Lemma~\ref{lem:descendant} therefore give a
contribution $O(\e^{-cH})$.  Restoring the other discarded cylinder unions,
and summing all polynomially many Fourier terms, gives the uniform bound
\begin{equation}\label{eq:delta-n}
 \delta_n:=L^{O_{r,D}(1)}
 \bigl(\e^{-cL^{1/2}}+\e^{-cH}+\rho^{cH}\bigr)=O_A(L^{-A}).
\end{equation}
Thus every nonzero Fourier term vanishes to the required order, while the
zero terms give the product on the right of
\eqref{eq:marked-factorization}.
\end{proof}

We also need fixed-radius blocks.  The recurrence
\eqref{eq:theta-cocycle} implies by induction that every torus character in
the full torus window can be reduced, on its digit cylinder, to a character
of the two central coordinates.  More precisely, fix $R\ge0$, an index
$j\ge R+1$ for the one-sided process, and integers
$c_{-R-1},\ldots,c_R$, and put
\[
 w_{j,R}=(a_{j-R+1},\ldots,a_{j+R}),
\]
with the word understood as empty when $R=0$.  There are integers
$A_{w_{j,R}},B_{w_{j,R}}$, depending only on this word, such that
\begin{equation}\label{eq:block-character-reduction}
 \sum_{t=-R-1}^{R}c_t\theta_{j+t}
 \equiv A_{w_{j,R}}\theta_j+B_{w_{j,R}}\theta_{j-1}\pmod1.
\end{equation}
Indeed, the forward substitutions up to $\theta_{j+R}$ use precisely
$a_{j+1},\ldots,a_{j+R}$, while the backward substitutions down to
$\theta_{j-R-1}$ use precisely $a_j,\ldots,a_{j-R+1}$.  On each fixed finite
collection of these words, $A_w,B_w$ range over a finite set.  In view of
\eqref{eq:theta}, the resulting phase is
$n(-1)^j(A_wq_j-B_wq_{j-1})\alpha$.

Let $R,K$ be fixed and call a cylinder--torus monomial a function of the form
\begin{equation}\label{eq:block-monomial}
 U_j(\alpha)=u(a_{j-R+1},\ldots,a_{j+R})
 \e^{2\pi i(r\theta_{j-1}+s\theta_j)},
\end{equation}
where $u$ is supported on finitely many words and $|r|+|s|\le K$.  Its
frequency is
\begin{equation}\label{eq:block-frequency}
 (-1)^jQ_j(r,s),\qquad Q_j(r,s)=sq_j-rq_{j-1}.
\end{equation}

\begin{proposition}[Two-block factorization]\label{prop:two-block}
For two fixed finite linear combinations $U,V$ of monomials
\eqref{eq:block-monomial}, all but $O_{U,V}(LH)$ pairs
$j<k$ in $\mathcal J_n$ satisfy
\begin{equation}\label{eq:two-block}
 \left|\int_0^1U_jV_k\,d\alpha
 -\int U\,d\widehat\mu_0\int V\,d\widehat\mu_0\right|
 \le C_{U,V}\bigl(\e^{-c_{\rm ld}L^{1/2}}
                    +\e^{-c_{\rm osc}H}
                    +\rho^{c_{\rm mix}H}\bigr),
\end{equation}
where $c_{\rm ld},c_{\rm osc},c_{\rm mix}>0$ may depend on $U,V$, and
$\widehat\mu_0$ is the stationary Gauss natural-extension measure times
Haar measure on $\T^2$.
\end{proposition}

\begin{proof}
It suffices to treat monomials; let $R$ be the maximum of their two radii.
If both torus modes vanish, ordinary
two-sided Gauss mixing proves \eqref{eq:two-block} unless the blocks overlap
or are separated by only $O(H)$ positions; these exceptional pairs number
$O(LH)$.

Choose $\kappa,\delta>0$ so that
\[
 \kappa+3\delta<80\lambda,
 \qquad \gamma:=80\lambda-\kappa-3\delta>0.
\]
Here $\kappa$ is used only for anti-concentration and $\delta$ only for the
denominator windows.  Put $c_{\rm anti}=2\log\varphi$, as supplied uniformly
by Lemma~\ref{lem:anti}, and write $c_{\rm ld}>0$
for the constant supplied by \eqref{eq:local-q-good} for this $\delta$.

Assume first that the later mode $(r_2,s_2)$ is nonzero.  Apply
Lemma~\ref{lem:anti} with $\eta=\e^{-\kappa H}$.  Since $k\ge200H$, outside a
complete depth-$k$ cylinder union of mass
\[
 O(\e^{-\kappa H}+\e^{-200c_{\rm anti}H})
 =O(\e^{-c_{\rm ac}H}),
 \qquad c_{\rm ac}:=\min(\kappa,200c_{\rm anti}),
\]
one has $|Q_k(r_2,s_2)|\ge\e^{-\kappa H}q_k$.  If
$k-j\ge C_{U,V,\kappa}H$, Fibonacci growth makes the earlier frequency at
most half this quantity, and hence
\[
 |Q|:=|(-1)^jQ_j(r_1,s_1)+(-1)^kQ_k(r_2,s_2)|
 \ge\tfrac12\e^{-\kappa H}q_k.
\]
The excluded separations again account for only $O_{U,V}(LH)$ pairs.  On
each depth-$(k+R)$ cylinder, the integer $Q$ and the complete amplitude are
fixed.  Put
\[
 t_-:=\left\lfloor\frac{m_n+k}{2}-40H\right\rfloor.
\]
Retain the complete depth-$k$ cylinders on which
$q_k\ge \e^{\lambda k-\delta H}$ and, below each such cylinder, the complete
depth-$t_-$ descendants on which
$q_{t_-}\le \e^{\lambda t_-+\delta H}$.  The discarded union has mass
$O(\e^{-c_{\rm ld}L^{1/2}})$.  On every retained descendant,
\begin{align*}
 \log\frac{q_{t_-}^2}{n|Q|}
 &\le 2\lambda t_-+2\delta H
       -L-\lambda k+(\kappa+\delta)H+O(1)\\
 &\le-(80\lambda-\kappa-3\delta)H+O(1)
 =-\gamma H+O(1).
\end{align*}
Thus, for all sufficiently large $n$,
$q_{t_-}^2\le\e^{-\gamma H/2}n|Q|$.  Lemma~\ref{lem:descendant}, with
$\varepsilon=\e^{-\gamma H/2}$, prefix depth $k+R$, and descendant depth
$t_-$, bounds the retained integral by $O(\e^{-\gamma H/2})$.  The summation
is by complete prefix-cylinder mass, so no cylinder-count factor is lost.
The stationary product is zero because the later Haar character is
nontrivial.

It remains to suppose that the earlier mode is nonzero and the later mode is
zero.  Write $Q_j=Q_j(r_1,s_1)$ and
\[
 t_0:=\frac{m_n+j}{2},\qquad
 t_-:=\lfloor t_0-40H\rfloor,
 \qquad t_+:=\lfloor t_0+40H\rfloor.
\]
Apply Lemma~\ref{lem:anti} with the same
$\eta=\e^{-\kappa H}$.  Outside a complete depth-$j$ cylinder union of mass
$O(\e^{-c_{\rm ac}H})$, one has
$|Q_j|\ge\e^{-\kappa H}q_j$.  Declare
$|k-t_0|\le100H$ exceptional.  This gives $O(H)$ values of $k$ per $j$,
hence $O(LH)$ pairs in total.

Retain throughout the complete depth-$j$ cylinder union on which
\[
 \e^{\lambda j-\delta H}\le q_j\le\e^{\lambda j+\delta H}.
\]
Its complement has mass $O(\e^{-c_{\rm ld}L^{1/2}})$, and on the retained
cylinders
\[
 \e^{-\kappa H}q_j\le |Q_j|\le C_{r_1,s_1}q_j.
\]

If $k<t_0-100H$, the whole two-block amplitude is measurable before depth
$t_-$.  Refine the retained depth-$j$ cylinders into complete
depth-$(k+R)$ prefixes, and below them keep the complete
depth-$t_-$ descendants with $q_{t_-}\le\e^{\lambda t_-+\delta H}$.  Their
complement has mass $O(\e^{-c_{\rm ld}L^{1/2}})$, while the same calculation
as above gives
\[
 \log\frac{q_{t_-}^2}{n|Q_j|}\le-\gamma H+O(1).
\]
Hence Lemma~\ref{lem:descendant}, with
$\varepsilon=\e^{-\gamma H/2}$, prefix depth $k+R$, and descendant depth
$t_-$, removes the phase with error
$O(\e^{-\gamma H/2})$.

If $k>t_0+100H$, also retain the complete depth-$t_+$ cylinders satisfying
$q_{t_+}\ge\e^{\lambda t_+-\delta H}$.  Their complement again has mass
$O(\e^{-c_{\rm ld}L^{1/2}})$.  Using the upper bound
$|Q_j|\le C_{r_1,s_1}q_j$ gives
\[
 \log\frac{q_{t_+}^2}{n|Q_j|}
 \ge(80\lambda-3\delta)H+O(1).
\]
Put $\gamma_+:=80\lambda-3\delta>0$.  For all sufficiently large $n$,
$q_{t_+}^2\ge\e^{\gamma_+H/2}n|Q_j|$, so
\eqref{eq:phase-freeze}, with $\varepsilon=\e^{-\gamma_+H/2}$, makes the
earlier phase constant up to $O(\e^{-\gamma_+H/2})$ on each retained
depth-$t_+$ cylinder.  Since the radii of $U,V$ are fixed, the later
zero-mode block starts at least $c_{\rm mix}H$ after $t_+$ for some
$c_{\rm mix}>0$ and all large $n$.  Conditional Gauss mixing, applied on
each cylinder and summed by cylinder mass, therefore replaces that block by
its stationary mean with error $O(\rho^{c_{\rm mix}H})$.  Extend this
replacement to the discarded depth-$t_+$ cylinders and restore them, at cost
$O(\e^{-c_{\rm ld}L^{1/2}})$; no future cutoff indicator then remains in the
earlier integral.  Finally retain the depth-$t_-$ descendants as above and
refine the remaining earlier amplitude to complete depth-$(j+R)$ prefixes.
Apply Lemma~\ref{lem:descendant}, with prefix depth $j+R$ and descendant
depth $t_-$, to the remaining earlier phase.

The stationary product is zero because the earlier Haar character is
nontrivial.  Taking
\[
 c_{\rm osc}:=\min(c_{\rm ac},\gamma/2,\gamma_+/2)
\]
and combining the complete-cylinder errors proves \eqref{eq:two-block}.
In particular, the denominator-deviation contribution is kept separate from
the anti-concentration and oscillatory exponents.
\end{proof}



\section{The principal Poisson and stable limits}

Let
\[
 Z_{n,j}=a_{j+1}W(\theta_j).
\]
If $A$ has the stationary Gauss digit law and $\Theta$ is independent and
uniform on $\T$, then
\begin{equation}\label{eq:stationary-tail}
 \Pp(AW(\Theta)>t)\sim\frac{c_0}{t},
 \qquad c_0=\frac1{12\log2}.
\end{equation}
Indeed,
\[
 \nu(A\ge m)=\frac1{\log2}\log\left(1+\frac1m\right)
 \sim\frac1{m\log2},
\]
and dominated convergence gives
\[
 \lim_{t\to\infty}t\Pp(AW(\Theta)>t)
 =\frac1{\log2}\int_0^1W(\theta)\,d\theta
 =\frac1{12\log2}.
\]

Define the marked point process
\begin{equation}\label{eq:Xi-def}
 \Xi_n:=\sum_{j\in\mathcal J_n}\delta_{(-1)^jZ_{n,j}/L}.
\end{equation}

\begin{proposition}[Poisson point process]\label{prop:PPP}
In the space of Radon point measures on $\R\setminus\{0\}$ equipped with
the vague topology,
\begin{equation}\label{eq:PPP}
 \Xi_n
 \ \Longrightarrow\
 \operatorname{PRM}(\Lambda),
 \qquad
 d\Lambda(x)=\frac1{2\pi^2}|x|^{-2}\,dx.
\end{equation}
\end{proposition}

\begin{proof}
We verify factorial moment measures against
$\varphi_1,\ldots,\varphi_r\in C_c^1(\R\setminus\{0\})$.
Fix an auxiliary exponent $A_0>0$, to be chosen sufficiently large.
Writing $X_{n,j}=(-1)^jZ_{n,j}/L$, the relevant ordered factorial sum is
\begin{equation}\label{eq:ordered-factorial-sum}
 \sum_{\substack{(j_1,\ldots,j_r)\in\mathcal J_n^r\\
                  j_1,\ldots,j_r\ \text{ pairwise distinct}}}
 \E\prod_{\ell=1}^r\varphi_\ell(X_{n,j_\ell}).
\end{equation}
Let $c>0$ be smaller than the distance of all supports from zero.  If
$\varphi_\ell((-1)^jaW(\theta)/L)\ne0$, then $a\ge8cL$.

First cut at $a\le L^D$, where $D>2$.  By
\eqref{eq:multi-digit-tail}, the total contribution of $r$-tuples for which
one digit exceeds $L^D$ and all the others exceed $8cL$ is
\[
 O\left(L^r\frac1{L^D L^{r-1}}\right)=O(L^{1-D}).
\]
For $a\le L^D$, the $C^1$ function
$f_{a,j}(\theta)=\varphi_\ell((-1)^jaW(\theta)/L)$ satisfies
$\|f_{a,j}'\|_\infty\le C_\varphi a/L\le C_\varphi L^{D-1}$.
Jackson's theorem \cite[Vol.~I, Chap.~III]{Zygmund59}, with degree
$M=L^{A_0+r+D+4}$, gives a trigonometric
polynomial whose uniform error is $O(L^{-A_0-r-4})$.  Its coefficient
$\ell^1$ norm is at most
$(2M+1)(\|\varphi_\ell\|_\infty+1)$ and hence is polynomial in $L$.
Increasing the complexity exponent in
Proposition~\ref{prop:marked-factorization}, we may apply that proposition;
after summation over $O(L^r)$ tuples, the approximation error is
$O(L^{-A_0-4})$.

Non-good time tuples are $O_r(L^{r-1}H)$ in number.  On each of them the
probability that all factors are nonzero is $O_r(L^{-r})$ by
\eqref{eq:multi-digit-tail}; their total contribution is $O(H/L)=o(1)$.
On good tuples, Proposition~\ref{prop:marked-factorization} factors the
expectation, and its error remains $o(1)$ after summing $O(L^r)$ tuples.
Equation \eqref{eq:stationary-tail} gives the vague convergence
\[
 L\Pp(AW(\Theta)/L\in dx)\Longrightarrow c_0x^{-2}\,dx
 \quad(x>0).
\]
To evaluate \eqref{eq:ordered-factorial-sum}, partition the ordered distinct
tuples according to the unique permutation $\sigma\in S_r$ for which
$j_{\sigma(1)}<\cdots<j_{\sigma(r)}$, and then according to the parity vector
$\varepsilon_\ell=(-1)^{j_\ell}\in\{-1,1\}$.  For every fixed
$\sigma$ and $\varepsilon=(\varepsilon_1,\ldots,\varepsilon_r)$,
\begin{equation}\label{eq:parity-tuple-count}
 \#\left\{(j_1,\ldots,j_r):
 j_{\sigma(1)}<\cdots<j_{\sigma(r)},\
 (-1)^{j_\ell}=\varepsilon_\ell\right\}
 =\frac{1+o(1)}{r!}\left(\frac{L}{2\lambda}\right)^r.
\end{equation}
Indeed, the parity-restricted lattice discrepancy in the actual bulk interval
is $O(L^{r-1})$, while replacing its length by $L/\lambda$ contributes
$O(HL^{r-1})=o(L^r)$.  Put
\[
 I_\ell^{\varepsilon}
 =\int_0^\infty\varphi_\ell(\varepsilon x)x^{-2}\,dx.
\]
For a fixed ordering and parity vector, factorization and
\eqref{eq:stationary-tail} show that the contribution tends to
\[
 \frac1{r!}\left(\frac{c_0}{2\lambda}\right)^r
 \prod_{\ell=1}^r I_\ell^{\varepsilon_\ell}.
\]
Summing first over the $r!$ temporal orderings and then over all parity
vectors yields
\[
 \prod_{\ell=1}^r
 \left\{\frac{c_0}{2\lambda}
 \int_{\R}\varphi_\ell(x)|x|^{-2}\,dx\right\}
 =\prod_{\ell=1}^r\int\varphi_\ell\,d\Lambda,
 \qquad \frac{c_0}{2\lambda}=\frac1{2\pi^2}.
\]
Thus the ordered factorial moment measures converge to those of the Poisson
random measure in \eqref{eq:PPP}.  If $K\Subset\R\setminus\{0\}$ and
$N_n(K)$ denotes the number of displayed points in $K$, the same iterated
conditional tail bound gives, for every $r\ge1$,
\[
 \sup_n\E (N_n(K))_r\le C_K^r,
\]
where $(x)_r=x(x-1)\cdots(x-r+1)$.  In particular all local factorial
moments are uniformly bounded, while the first moment supplies local
tightness.  Moreover, for $0\le u\le1$,
\[
 \sup_n\E(1+u)^{N_n(K)}
 =\sup_n\sum_{r\ge0}\frac{u^r}{r!}\E(N_n(K))_r
 \le \exp(C_Ku).
\]
Thus the local counts have uniformly bounded exponential moments.  On
finitely many disjoint compact continuity sets, the preceding calculation
gives convergence of all mixed factorial moments to those of independent
Poisson variables, and the exponential bound makes those moments
distribution determining.  Local tightness and an exhaustion by compact
continuity sets now identify every subsequential limit as the Poisson random
measure with intensity $\Lambda$; this is the standard factorial-moment
criterion \cite[Chap.~4]{Kallenberg17}.
\end{proof}

We next control the compensated small jumps.  Choose
$\chi\in C^1([0,\infty))$ with $0\le\chi\le1$, $\chi=1$ on $[0,1]$, and
$\chi=0$ on $[2,\infty)$.  For $0<\varepsilon<1$, set
\begin{equation}\label{eq:small-jump-U}
 U_{n,j}^{(\varepsilon)}=(-1)^jZ_{n,j}
 \chi\!\left(\frac{Z_{n,j}}{\varepsilon L}\right).
\end{equation}

\begin{lemma}[Small jumps]\label{lem:small-jumps}
For a constant $C$ independent of $\varepsilon$,
\begin{equation}\label{eq:small-jump-var}
 \limsup_{n\to\infty}\frac1{L^2}
 \Var\left(\sum_{j\in\mathcal J_n}U_{n,j}^{(\varepsilon)}\right)
 \le C\varepsilon.
\end{equation}
\end{lemma}

\begin{proof}
Since $W\le1/8$, Lemma~\ref{lem:gauss-estimates}(ii) gives, uniformly in
$j,n$,
\[
 \Pp(Z_{n,j}>t)\le\frac{C}{1+t}.
\]
For $j\ne k$, \eqref{eq:multi-digit-tail} similarly gives
\[
 \Pp(Z_{n,j}>s,Z_{n,k}>t)
 \le\frac{C}{(1+s)(1+t)}.
\]
Layer-cake integration, with $M=2\varepsilon L$, yields
\begin{equation}\label{eq:small-moments}
 \E|U_{n,j}^{(\varepsilon)}|^2\le C\varepsilon L,
 \qquad
 \E|U_{n,j}^{(\varepsilon)}U_{n,k}^{(\varepsilon)}|
 \le C\log^2(2+L).
\end{equation}
The diagonal contribution is $O(\varepsilon L^2)$.  Overlapping,
$H$-near, or resonance-near pairs number $O(LH)$, and their total
contribution is $O(LH\log^2L)=o(L^2)$.

For the remaining pairs, cut each digit at $a\le A_L=L^D$.  Since
$|U_{n,j}^{(\varepsilon)}|\le2\varepsilon L$,
\[
 \|U_{n,j}^{(\varepsilon)}\ind_{\{a_{j+1}>A_L\}}\|_2
 \le C\varepsilon L A_L^{-1/2}.
\]
After centering, $\|X-\E X\|_2\le2\|X\|_2$; hence, after division by $L$,
Minkowski's inequality shows that the sum of these tails is
$O(\varepsilon L^{1-D/2})=o(1)$ if $D>2$.  Fix an auxiliary exponent
$A_0>0$, to be chosen sufficiently large.  The map
$z\mapsto z\chi(z/(\varepsilon L))$ has a Lipschitz constant bounded
independently of $\varepsilon$ and $L$ on
$[0,\infty)$, and $W$ is $1/2$-Lipschitz on $\T$.  Hence, on $a\le A_L$,
the truncated summand has Lipschitz norm at most $Ca\le CL^D$, uniformly in
$0<\varepsilon<1$.  Jackson approximation for Lipschitz functions, with
uniform error $L^{-A_0-4}$ and polynomial degree and coefficient norm, applies
as in the proof of Proposition~\ref{prop:PPP}.
Centering doubles the approximation error but does not change its order.
Proposition~\ref{prop:marked-factorization},
with $A_0$ chosen so large that $L^2O_{A_0}(L^{-A_0})=o(1)$, makes the sum of all
remaining covariances $o(L^2)$.  This proves
\eqref{eq:small-jump-var}.
\end{proof}

\begin{corollary}[Principal Cauchy law]\label{cor:principal-cauchy}
There are deterministic numbers $b_n^{(0)}$ such that
\begin{equation}\label{eq:principal-cauchy}
 \frac1L\sum_{j\in\mathcal J_n}(-1)^jZ_{n,j}-b_n^{(0)}
 \Longrightarrow \operatorname{Cauchy}\!\left(0,\frac1{2\pi}\right).
\end{equation}
\end{corollary}

\begin{proof}
For fixed $\varepsilon$, split $x=x\chi(|x|/\varepsilon)
+x(1-\chi(|x|/\varepsilon))$.  Proposition~\ref{prop:PPP} and the
continuous mapping theorem give convergence of the second, large-jump part;
the digit-tail bound gives the uniform finite-$n$ estimate
\[
 \sup_n\E\Xi_n\{x:|x|>R\}\le C/R,
\]
so the tails $|x|>R$ may be removed before applying continuous mapping.
Center the small part by
\[
 b_{n,\varepsilon}=\frac1L\sum_{j\in\mathcal J_n}
 \E U_{n,j}^{(\varepsilon)}.
\]
Put
\[
 R_{n,\varepsilon}
 =\frac1L\sum_{j\in\mathcal J_n}U_{n,j}^{(\varepsilon)}
   -b_{n,\varepsilon},
 \qquad
 Y_{n,\varepsilon}
 =\int x\bigl(1-\chi(|x|/\varepsilon)\bigr)\,d\Xi_n(x).
\]
Then the principal sum centered by $b_{n,\varepsilon}$ is exactly
$Y_{n,\varepsilon}+R_{n,\varepsilon}$, and Lemma~\ref{lem:small-jumps}
and Chebyshev's inequality give, for every $\delta>0$,
\begin{equation}\label{eq:small-jump-prob}
 \limsup_{n\to\infty}
 \Pp(|R_{n,\varepsilon}|>\delta)
 \le \frac{C\varepsilon}{\delta^2}.
\end{equation}
Thus the centered small part becomes negligible as $\varepsilon\downarrow0$.
The transition shell
$\varepsilon<|x|<2\varepsilon$ is included in the same continuous
truncation; its limiting compensated variance is $O(\varepsilon)$.

For fixed $\varepsilon$, Proposition~\ref{prop:PPP} and the preceding tail
removal give $Y_{n,\varepsilon}\Longrightarrow Y_\varepsilon$.  For the
smooth cutoff used here, the characteristic exponent of $Y_\varepsilon$ is,
first on $|x|\le R$,
\[
 \int_{|x|\le R}
 \left(\exp\{itx(1-\chi(|x|/\varepsilon))\}-1\right)d\Lambda(x).
\]
Letting $R\to\infty$ is legitimate because the expected number of points
outside $[-R,R]$ is $O(R^{-1})$.  Symmetry cancels the imaginary part, and
then dominated convergence as $\varepsilon\downarrow0$ gives
\[
 \int_{\R}(\cos(tx)-1)\frac{dx}{2\pi^2x^2}
 =-\frac{|t|}{2\pi}.
\]
Hence $Y_\varepsilon\Longrightarrow
\operatorname{Cauchy}(0,1/(2\pi))$ as $\varepsilon\downarrow0$.
The convergence-together criterion applied to this limit and
\eqref{eq:small-jump-prob} now proves the claim.  Explicitly, choose
$\varepsilon_m\downarrow0$ so that $\varepsilon_m m^2\to0$ and the
Prokhorov distance from the law of $Y_{\varepsilon_m}$ to the target law is
at most $1/m$.  Then choose $N_m\uparrow\infty$ so that, for $n\ge N_m$,
the Prokhorov distance between the laws of $Y_{n,\varepsilon_m}$ and
$Y_{\varepsilon_m}$ is at most $1/m$ and
$\Pp(|R_{n,\varepsilon_m}|>1/m)\le2C\varepsilon_m m^2$.  Set
$\varepsilon_n=\varepsilon_m$ for $N_m\le n<N_{m+1}$, and define it
arbitrarily for $n<N_1$.  With
$b_n^{(0)}:=b_{n,\varepsilon_n}$ this gives
\eqref{eq:principal-cauchy}.
\end{proof}





\section{The carry graph and the bounded remainder}

Let $(\widehat\Omega,\widehat\nu,\sigma)$ be the two-sided natural extension
of the Gauss system.  Concretely, up to null boundary sets, we use
\[
 \widehat\Omega=(0,1)^2,
 \qquad
 d\widehat\nu(x,y)=\frac{dx\,dy}{\log2\,(1+xy)^2},
 \qquad
 \sigma(x,y)=\left(Tx,\frac1{a_1(x)+y}\right),
\]
and write $x(\omega)=[0;a_1(\omega),a_2(\omega),\ldots]$.
Whenever order, ceiling, or fractional part is used for a torus coordinate,
we identify $\T$ with its canonical representatives in $[0,1)$.
Define the torus cocycle
\begin{equation}\label{eq:S-cocycle}
 \mathcal S(\omega,r,s)=
 (\sigma\omega,s,r-a_1(\omega)s\pmod1).
\end{equation}
The fibre matrix
\[
 A_a=\begin{pmatrix}0&1\\1&-a\end{pmatrix}
\]
lies in $\operatorname{GL}_2(\Z)$ and preserves Haar measure on $\T^2$.
The displayed density $\widehat\nu$ is $\sigma$-invariant by branchwise
change of variables; hence
$\widehat\mu_0=\widehat\nu\otimes m_{\T^2}$ is invariant.

\begin{lemma}[Endpoint-controlled continuant recurrence]
\label{lem:endpoint-recurrence}
Let $a_i\ge1$ and let an integer sequence with $(z_0,z_1)\ne(0,0)$ satisfy
\[
 z_{i+2}=a_{i+1}z_{i+1}+z_i\qquad(0\le i<m).
\]
If
\[
 \max(|z_0|,|z_1|,|z_m|,|z_{m+1}|)\le K,
\]
then $m\le C\log(2K)$ for an absolute constant $C$.
\end{lemma}

\begin{proof}
Put $S_i=|z_i|+|z_{i+1}|$ for $0\le i\le m$, and call $i$ aligned if
$z_i z_{i+1}\ge0$.  Since the recurrence is invertible, no adjacent pair is
$(0,0)$, and hence $S_i\ge1$.  Alignment propagates forwards,
because
\[
 z_{i+1}z_{i+2}
 =a_{i+1}z_{i+1}^2+z_i z_{i+1}\ge0.
\]
Applying the recurrence twice at an aligned index gives the first inequality
below.  If $i+1$ and $i+2$ are both unaligned, solving the same two steps
backwards gives the second:
\[
 S_{i+2}\ge2S_i,\qquad S_i\ge2S_{i+2}\qquad(i+2\le m),
\]
with the respective alignment hypotheses just stated.

Let $t$ be the first aligned index in $\{0,\ldots,m\}$, with $t=m+1$ if
there is none.  Set $k_-=\lfloor(t-1)/2\rfloor$ for $t\ge1$ and $k_-=0$
for $t=0$; set $k_+=\lfloor(m-t)/2\rfloor$ for $t\le m$ and $k_+=0$
otherwise.  Iteration on the unaligned prefix and aligned suffix gives
\[
 m\le2k_-+2k_++3,
 \qquad 2^{k_-}\le S_0,
 \qquad 2^{k_+}\le S_m.
\]
Here the possible final unpaired step is harmless because $S_i$ is
nondecreasing on the aligned suffix.  Since $S_0,S_m\le2K$ and the nonzero
integer initial pair forces $K\ge1$, taking logarithms now gives
\[
 m\le\frac{4\log(2K)}{\log2}+3
 \le\frac7{\log2}\log(2K).
\]
This proves the claim, for example with $C=7/\log2$.
\end{proof}

\begin{lemma}[Gauss--torus mixing]\label{lem:torus-mixing}
The system $(\widehat\Omega\times\T^2,\widehat\mu_0,\mathcal S)$ is mixing.
\end{lemma}

\begin{proof}
It is enough to use finite digit-cylinder functions times torus characters.
For characters $k,\ell\in\Z^2$, the fibre integral in an $m$-step
correlation vanishes unless
\begin{equation}\label{eq:character-resonance}
 k+(A_{a_m}\cdots A_{a_1})^{\mathrm T}\ell=0.
\end{equation}
Indeed, reverse the matrix order forced by the transpose and, for
$0\le i\le m$, put
\[
 w_i=A_{a_{m-i+1}}^{\mathrm T}\cdots A_{a_m}^{\mathrm T}\ell
     =(\xi_i,\eta_i)^{\mathrm T},
 \qquad w_0=\ell.
\]
Since $A_a^{\mathrm T}=A_a$, for $0\le i<m$,
\[
 w_{i+1}=A_{a_{m-i}}^{\mathrm T}w_i
          =(\eta_i,\xi_i-a_{m-i}\eta_i)^{\mathrm T}.
\]
Thus $\xi_{i+1}=\eta_i$; on defining $\xi_{m+1}:=\eta_m$, we have
\[
 \xi_{i+2}=\xi_i-a_{m-i}\xi_{i+1}
 \qquad(0\le i<m).
\]
Consequently $z_i=(-1)^i\xi_i$ satisfies
$z_{i+2}=a_{m-i}z_{i+1}+z_i$, which is the recurrence of
Lemma~\ref{lem:endpoint-recurrence} after relabelling the positive
coefficients.  Moreover,
$w_m=(A_{a_m}\cdots A_{a_1})^{\mathrm T}\ell=-k$ under
\eqref{eq:character-resonance}; hence $(z_0,z_1)$ is, up to signs, the
coordinate pair of $\ell$, and $(z_m,z_{m+1})$ is the coordinate pair of
$k$.  Since the matrix product in \eqref{eq:character-resonance} is
invertible, $\ell=0$ if and only if $k=0$.  In the nonzero case
$(z_0,z_1)\ne(0,0)$, and both endpoint pairs are bounded in terms of the
fixed vectors $k$ and $\ell$, so Lemma~\ref{lem:endpoint-recurrence} applies.
Thus the resonance relation is impossible for all sufficiently large $m$.
The zero modes reduce to mixing of the Gauss natural extension; this follows
from Lemma~\ref{lem:gauss-estimates}(i), since the natural extension of a
mixing probability-preserving endomorphism is mixing.
Density of cylinder functions times characters in $L^2$ completes the proof.
\end{proof}

We need to transfer fixed windows of the actual curve process to this
stationary system.  The complete quotient is an infinite-future coordinate,
so we record the required full-state form rather than silently treating it as
a finite digit function.

For a fixed $R\ge0$, give
\begin{equation}\label{eq:window-space}
 \mathcal X_R
 =\N^{2R+1}\times[0,1]^{2R+1}\times\T^{2R+2}
\end{equation}
the product topology, with $\N$ discrete, and define the actual window
\begin{equation}\label{eq:actual-window}
 \mathcal W_{n,j}^{(R)}=
 \left((a_{j+t+1})_{t=-R}^{R},
       (x_{j+t})_{t=-R}^{R},
       (\theta_{j+t})_{t=-R-1}^{R}\right)\in\mathcal X_R.
\end{equation}
The stationary window is defined from the two-sided Gauss--torus system in
the same way, and its law is denoted by $\mu_R$.  Both laws are supported on
the corresponding orbit-consistent Borel subset of $\mathcal X_R$.

For $R'\ge R$, let
$\pi_{R',R}:\mathcal X_{R'}\to\mathcal X_R$ be the coordinate projection
that retains the digit and real coordinates with $-R\le t\le R$ and the
torus coordinates with $-R-1\le t\le R$.  Then
$(\pi_{R',R})_*\mu_{R'}=\mu_R$, and, whenever both actual windows are
defined,
$\pi_{R',R}(\mathcal W_{n,j}^{(R')})=\mathcal W_{n,j}^{(R)}$.

\begin{lemma}[Uniform full-state transfer]\label{lem:full-state}
Let $\eta_{n,j}^{(R)}$ be the law of \eqref{eq:actual-window} for Lebesgue
$\alpha$, and let $\mu_R$ be the stationary window law.  If
$G:\mathcal X_R\to\mathbb C$ is bounded and Borel measurable, and its
discontinuity set has
$\mu_R$-measure zero, then
\begin{equation}\label{eq:full-state}
 \sup_{j\in\mathcal J_n}
 \left|\int G\,d\eta_{n,j}^{(R)}-\int G\,d\mu_R\right|\longrightarrow0.
\end{equation}
\end{lemma}

\begin{proof}
First take a digit-cylinder amplitude supported on finitely many local words,
times a torus character
$\exp(2\pi i\sum_{t=-R-1}^{R}c_t\theta_{j+t})$, and let $d$ be the largest
initial digit depth on which the chosen test depends.  For the raw digit block
in \eqref{eq:actual-window}, $d=j+R+1$; after replacing complete quotients by
a fixed number $M$ of future digits below, one may have $d\le j+R+M$.
Partition into complete depth-$d$ cylinders over the finitely many supported
local words.  By \eqref{eq:block-character-reduction}, on each such cylinder
the coefficients $A_w,B_w$, the integer
$Q_j=A_wq_j-B_wq_{j-1}$, and the whole block amplitude are fixed.  If $A_w=B_w=0$, uniform Gauss mixing gives the stationary mean.
Otherwise choose $\kappa,\delta>0$ with
$\kappa+3\delta<80\lambda$ and apply Lemma~\ref{lem:anti} with
$\eta=\e^{-\kappa H}$.  The bound there is uniform in the nonzero pair
$(A_w,B_w)$, so after summing over the finitely many words a complete cylinder
union of mass
$O(\e^{-\kappa H}+F_{j+1}^{-2})=O(\e^{-cH})$ is removed, and on its complement
$|Q_j|\ge\e^{-\kappa H}q_j$.

Put $t=\lfloor(m_n+j)/2-40H\rfloor$.  Because
$j\in\mathcal J_n$ and $d-j=O_{R,M}(1)$ for each fixed cylinder test, one has
$t>d$ for all sufficiently large $n$, uniformly in $j$.  Retain the coarser
complete depth-$j$ cylinder union on which
$q_j\ge\e^{\lambda j-\delta H}$ and, below every retained depth-$d$ prefix,
the complete depth-$t$ descendants on which
$q_t\le\e^{\lambda t+\delta H}$.  The discarded mass is
$O(\e^{-c_\delta L^{1/2}})$ by \eqref{eq:local-q-good}.  On every retained
descendant,
\[
 \log\frac{q_t^2}{n|Q_j|}
 \le-(80\lambda-\kappa-3\delta)H+O(1),
\]
and hence $q_t^2\le\e^{-c'_{R,G}H}n|Q_j|$ for some $c'_{R,G}>0$ and all
large $n$.  Lemma~\ref{lem:descendant}, with prefix depth $d$, therefore
makes every nonzero character
$O(\e^{-c'_{R,G}H}+\e^{-c_\delta L^{1/2}})$.  This proves the required one-block convergence uniformly in $j$ for finite
sums of cylinder--character tests with finitely supported local digit
amplitudes.

General finite-window tests are reduced to this class by digit truncation.
In particular, first restrict the finitely many digits required by
\eqref{eq:block-character-reduction} and by the digit-cylinder amplitude to
$a\le K$; Lemma~\ref{lem:gauss-estimates}(ii) and a union bound make the
omitted mass $O_R(1/K)$, uniformly for the actual and stationary window laws.
Continuous functions of the complete quotients are then included by
replacing, simultaneously for every $|t|\le R$, the coordinate
$x_{j+t}=[0;a_{j+t+1},a_{j+t+2},\ldots]$ by its first $M$ future digits.  The
maximum coordinate error is $O_R(F_M^{-2})$; hence, for continuous $G$, the
test-function error is at most $\omega_G(C_RF_M^{-2})$.  Restrict every digit
appearing in these finite continued fractions to $a\le K$ as well.  The total
discarded mass is $O_R(M/K)$, uniformly for both window laws.  These truncated
tests are covered by the preceding argument with $d\le j+R+M$.  For fixed
$M,K$ only finitely many words remain; first let $n\to\infty$, then
$K\to\infty$, and finally $M\to\infty$.

For $K\ge1$, let $\mathcal X_{R,K}\subset\mathcal X_R$ be the set on which
all digit coordinates are at most $K$; it is compact, and these sets exhaust
$\mathcal X_R$.  The unital self-conjugate algebra generated by finite digit-cylinder
indicators, continuous functions of the real coordinates, and torus
characters separates points on each $\mathcal X_{R,K}$.
Stone--Weierstrass therefore makes it uniformly dense there;
the same algebra is convergence determining and its real span is dense in
$L^2(\mu_R)$.  The family $\eta_{n,j}^{(R)}$ is uniformly tight: real and
torus coordinates are compact, and Lemma~\ref{lem:gauss-estimates}(ii)
controls the finitely many digit coordinates.  If \eqref{eq:full-state}
failed, tightness would give a weakly convergent subsequence, say with limit
$\eta$.  Convergence on the determining algebra forces $\eta=\mu_R$, and the
Portmanteau theorem applied to the bounded $\mu_R$-almost-everywhere
continuous function $G$ then contradicts the assumed failure.
\end{proof}

For $z=(x,r,s)$ and $d\in\N_0$, define
\begin{equation}\label{eq:carry-map}
 \phi_z(d)=\lceil x(d+r)-s\rceil.
\end{equation}
Equation \eqref{eq:u-carry} shows that the actual carries satisfy
$d_{j+1}=\phi_{(x_j,\theta_{j-1},\theta_j)}(d_j)$.
If $e=d+r$, the updated real carry obeys
\begin{equation}\label{eq:carry-contraction}
 xe\le e'<xe+1.
\end{equation}
Combining this with \eqref{eq:fibonacci-product}, all actual trajectories and
all pullback trajectories started from zero remain below an absolute integer
$D$.

Set
\begin{equation}\label{eq:reset-set}
 \mathcal R=\left\{(x,r,s):
 x<\frac1{4(D+1)},\quad\frac12<s<\frac34\right\}.
\end{equation}
For $0\le d\le D$ and $z\in\mathcal R$,
\[
 -\frac34<x(d+r)-s<-\frac14,
\]
so $\phi_z(d)=0$.  The set $\mathcal R$ has positive
$\widehat\mu_0$-measure.  By Lemma~\ref{lem:torus-mixing}, almost every
two-sided orbit visits it infinitely often in the past.  Start at time $-N$
with carry zero and iterate to time zero.  After the most recent reset the
answer no longer depends on $N$, so it stabilizes almost surely to a
measurable function $d_*$ satisfying
\begin{equation}\label{eq:carry-graph}
 d_*(\mathcal Sz)=\phi_z(d_*(z)).
\end{equation}
Thus the stationary carry extension is a measurable invariant graph and is
isomorphic to the mixing base--torus system.

For $R\ge1$, define $d_j^{(R)}$ by setting the carry equal to zero at time
$j-R$ and iterating $R$ times.  Let
\[
 p_R=\widehat\mu_0\left(
 \bigcap_{t=0}^{R-1}\mathcal S^{-t}\mathcal R^c\right).
\]
Ergodicity and $\widehat\mu_0(\mathcal R)>0$ imply $p_R\to0$.  Whenever the
actual orbit has a reset during $[j-R,j)$, $d_j=d_j^{(R)}$.  The no-reset
indicator has stationary-null boundary, so Lemma~\ref{lem:full-state} gives
\begin{equation}\label{eq:carry-coupling-prob}
 \sup_{j\in\mathcal J_n}\Pp(d_j\ne d_j^{(R)})\le p_R+o_n(1).
\end{equation}

We now prove the weak law for the remainder in
\eqref{eq:principal-split}.

\begin{proposition}[Bounded-remainder weak law]\label{prop:remainder}
\begin{equation}\label{eq:remainder-WLLN}
 \frac1L\sum_{j\in\mathcal J_n}(-1)^j
 (B_j-\E B_j)\longrightarrow0
 \quad\text{in probability}.
\end{equation}
\end{proposition}

\begin{proof}
Let $B_j^{(R)}$ be the bounded remainder obtained by replacing $d_j$ with
$d_j^{(R)}$.  Choose a bounded Borel function
$B^{(R)}:\mathcal X_R\to\R$ whose restriction to the orbit-consistent support
represents this remainder, so that
$B_j^{(R)}=B^{(R)}(\mathcal W_{n,j}^{(R)})$.  For fixed $R$, this function is
bounded and $\mu_R$-almost-everywhere continuous: its discontinuities lie on
countably many Gauss boundaries and finitely many ceiling or fractional-part
hypersurfaces, all of stationary measure zero.

\medskip
\noindent\textit{Density bridge.}
Fix $R$ and $\varepsilon>0$.  The $L^2(\mu_R)$-density established after
Lemma~\ref{lem:full-state} first gives a finite sum
\[
 G(\mathbf a,\mathbf x,\boldsymbol\vartheta)
 =\sum_{\ell=1}^N D_\ell(\mathbf a)g_\ell(\mathbf x)
   \exp\!\left(2\pi i\sum_{t=-R-1}^{R}
                    c_{\ell,t}\vartheta_t\right),
\]
where each $D_\ell$ is a finite digit-cylinder function, each $g_\ell$ is
continuous on $[0,1]^{2R+1}$, and each $c_{\ell,t}$ is an integer, such that
\[
 \|B^{(R)}-G\|_{L^2(\mu_R)}<\varepsilon/4.
\]
The torus range here is the full range $-R-1\le t\le R$ from
\eqref{eq:actual-window}.

For an integer $M\ge1$, put $R_M=R+M$ and lift both functions to
$\mathcal X_{R_M}$ by $\pi_{R_M,R}$.  Since
$(\pi_{R_M,R})_*\mu_{R_M}=\mu_R$, the first $L^2$ error remains less than
$\varepsilon/4$ after this lift.  On an orbit-consistent window replace,
for every $-R\le t\le R$,
\[
 x_{j+t}=[0;a_{j+t+1},a_{j+t+2},\ldots]
 \quad\hbox{by}\quad
 x_{j+t}^{[M]}=[0;a_{j+t+1},\ldots,a_{j+t+M}].
\]
The continued-fraction contraction estimate gives
$\max_{|t|\le R}|x_{j+t}-x_{j+t}^{[M]}|=O_R(F_M^{-2})$.  Since only finitely
many continuous factors $g_\ell$ occur, their uniform continuity permits $M$
to be chosen so that the resulting function $G_M$ satisfies
\[
 \|G\circ\pi_{R_M,R}-G_M\|_{L^2(\mu_{R_M})}<\varepsilon/4.
\]

Now let $E_{M,K}$ be the event that every digit in the full
radius-$R_M$ monomial word
\[
 a_{j-R_M+1},\ldots,a_{j+R_M}
\]
is at most $K$.  This truncates all digits introduced by the finite-future
replacement and also the $M$ left-padding digits needed to regard the result
as a radius-$R_M$ monomial.  By Lemma~\ref{lem:gauss-estimates}(ii),
stationarity, and a union bound,
\[
 \mu_{R_M}(E_{M,K}^c)=O_R(M/K).
\]
Since $G_M$ is bounded, this measure estimate gives the required $L^2$
estimate explicitly:
\[
 \|G_M-G_M\ind_{E_{M,K}}\|_{L^2(\mu_{R_M})}^2
 \le \|G_M\|_\infty^2\mu_{R_M}(E_{M,K}^c)
 =O_{R,G}(M/K).
\]
Choose $K$ so large that the corresponding $L^2$ norm is less than
$\varepsilon/4$.  For fixed $M,K$, only finitely many full digit words remain.
For each summand $\ell$ and each such word $\widetilde w$, identity
\eqref{eq:block-character-reduction}, applied to the central subword
$w=(a_{j-R+1},\ldots,a_{j+R})$, gives
\[
 \sum_{t=-R-1}^{R}c_{\ell,t}\theta_{j+t}
 \equiv A_{\ell,w}\theta_j+B_{\ell,w}\theta_{j-1}\pmod1.
\]
Thus $G_M\ind_{E_{M,K}}$ is exactly a finite linear combination of the
monomials \eqref{eq:block-monomial} at radius $R_M=R+M$: its digit amplitude
is supported on the finitely many words $\widetilde w$, and its central
Fourier modes are $(r,s)=(B_{\ell,w},A_{\ell,w})$.  These modes have finite
support because both the original torus-character list and the retained word
list are finite.
The radius grows with $M$ because the approximation to $x_{j+R}$ uses digits
through $a_{j+R+M}$; the additional left coordinates in the symmetric
radius-$R_M$ word are harmless and were truncated above.

Combining the three errors proves that, for every $\varepsilon>0$, there are
$M,K$ and a finite linear combination $P$ of monomials at radius $R+M$ such
that
\[
 \|B^{(R)}\circ\pi_{R+M,R}-P\|_{L^2(\mu_{R+M})}<\varepsilon.
\]
Because $B^{(R)}\circ\pi_{R+M,R}$ is real-valued,
\[
 |B^{(R)}\circ\pi_{R+M,R}-\operatorname{Re}P|
 \le |B^{(R)}\circ\pi_{R+M,R}-P|
\]
pointwise.  Adjoining the conjugate monomials shows that
$\operatorname{Re}P$ is again a finite linear combination of monomials and is
real-valued.  Taking a sequence of accuracies decreasing
to zero and choosing the truncation lengths increasingly, we let $M$ range
over that increasing sequence.  We thereby obtain real-valued $P_{R,M}$ at
radii $R_M=R+M$ such that
\begin{equation}\label{eq:remainder-L2-approx}
 \|B^{(R)}\circ\pi_{R_M,R}-P_{R,M}\|_{L^2(\mu_{R_M})}
 =:\delta_{R,M}\longrightarrow0.
\end{equation}

For each fixed $R,M$, apply Lemma~\ref{lem:full-state} with window radius
$R_M$ to the bounded, $\mu_{R_M}$-almost-everywhere continuous function
$|B^{(R)}\circ\pi_{R_M,R}-P_{R,M}|^2$ to obtain
\begin{equation}\label{eq:actual-L2-transfer}
 \sup_{j\in\mathcal J_n}
 \E|B_j^{(R)}-P_{R,M,j}|^2
 =\delta_{R,M}^2+o_n(1).
\end{equation}
For fixed $R,M$, every $P_{R,M}$ is, by the bridge, a finite linear
combination of the monomials \eqref{eq:block-monomial}; its radius, digit-word
support, and Fourier support are all fixed before $n\to\infty$.  Moreover,
its integral against $\mu_{R_M}$ is its integral against $\widehat\mu_0$ under
the stationary window map.  Hence Proposition~\ref{prop:two-block} applies to
every covariance term in
$\sum_{j\in\mathcal J_n}(-1)^j(P_{R,M,j}-\E P_{R,M,j})$.  Its variance is
$o(L^2)$: the diagonal contributes $O(L)$, the $O(LH)$ exceptional pairs
contribute $O(LH)=o(L^2)$, and all other pairs contribute
\[
 O_{R,M}\!\left(L^2\bigl(\e^{-c_{\rm ld}L^{1/2}}
                  +\e^{-c_{\rm osc}H}+\rho^{c_{\rm mix}H}\bigr)\right)
 =o(1).
\]
Lemma~\ref{lem:full-state} gives
$\sup_j|\E P_{R,M,j}-\int P_{R,M}\,d\mu_{R_M}|=o(1)$, so replacing the
stationary means in Proposition~\ref{prop:two-block} by the actual means adds
only $o(L^2)$ to the covariance sum.

Write $\widetilde X=X-\E X$.  Minkowski's inequality and
\eqref{eq:actual-L2-transfer} give
\begin{align}\label{eq:Minkowski-polynomial}
 &\left\|\frac1L\sum_{j\in\mathcal J_n}(-1)^j
 (\widetilde B_j^{(R)}-\widetilde P_{R,M,j})\right\|_2\\
 &\qquad\le\frac2L\sum_{j\in\mathcal J_n}
 \|B_j^{(R)}-P_{R,M,j}\|_2
 \le C\delta_{R,M}+o_n(1).\notag
\end{align}
First let $n\to\infty$, then $M\to\infty$.  This proves the analogue of
\eqref{eq:remainder-WLLN} with $B_j^{(R)}$ for every fixed $R$.

Finally, boundedness and \eqref{eq:carry-coupling-prob} imply
\[
 \sup_{j\in\mathcal J_n}\|B_j-B_j^{(R)}\|_2
 \le C\sqrt{p_R+o_n(1)}.
\]
A second centered Minkowski estimate yields
\begin{equation}\label{eq:Minkowski-reset}
 \limsup_{n\to\infty}
 \left\|\frac1L\sum_{j\in\mathcal J_n}(-1)^j
 (\widetilde B_j-\widetilde B_j^{(R)})\right\|_2
 \le C\sqrt{p_R}.
\end{equation}
Letting $R\to\infty$ completes the proof.  The order of limits is
$n\to\infty$, then $M\to\infty$, then $R\to\infty$.
\end{proof}




\section{Stopping, centering, and completion of the proof}

The continuous height has the exact form
\begin{equation}\label{eq:Cq}
 C_j=\frac{n}{q_j+q_{j-1}x_j},
 \qquad N_j=C_j-E_j,\qquad 0\le E_j\le E_*.
\end{equation}

\begin{lemma}[Stopping time and end terms]\label{lem:stopping}
With $H=L^{3/4}$,
\begin{equation}\label{eq:tau}
 \tau_n=\frac L\lambda+O_{\Pp}(H).
\end{equation}
Moreover,
\begin{equation}\label{eq:end-terms}
 \frac1L\left[
 S_n(\alpha)-\sum_{j\in\mathcal J_n}(-1)^j\Phi(x_j,u_j)
 \right]\longrightarrow0
 \quad\text{in probability}.
\end{equation}
\end{lemma}

\begin{proof}
Since $q_j\le q_j+q_{j-1}x_j<2q_j$, \eqref{eq:Cq} gives a deterministic
bracketing.  If $q_j\le n/(2(E_*+1))$, then $C_j\ge E_*+1$ and hence
$N_j\ge1$.  If $q_j>n$, then $C_j<1$ and, because $N_j$ is a nonnegative
integer, $N_j=0$.  Thus $\tau_n$ lies between the hitting times at which
$q_j$ crosses these two fixed multiples of $n$.  Applying
\eqref{eq:q-large-deviation} at $L/\lambda\pm CH$ to both thresholds, and
then choosing $C$ large, proves \eqref{eq:tau}.

Also,
\begin{equation}\label{eq:Phi-tail-bound}
 |\Phi(x,u)|\le\frac1{8x}+\frac12
 \le\frac{a_1(x)}8+\frac58.
\end{equation}
Lemma~\ref{lem:gauss-estimates}(ii) therefore gives
\begin{equation}\label{eq:Phi-tail}
 \sup_{j,n}\Pp(|\Phi(x_j,u_j)|>t)\le\frac{C}{1+t}.
\end{equation}
Let $I_n=\{0,\ldots,\tau_n-1\}$.  On the event
$|\tau_n-m_n|\le CH$, the symmetric difference
$I_n\mathbin{\triangle}\mathcal J_n$ is contained in the deterministic union
\[
 \{0\le j<200H+1\}
 \ \cup\ 
 \{m_n-(C+200)H-2\le j\le m_n+CH+2\},
\]
which has $O_C(H)$ positions.  Moreover,
\[
 \left|S_n(\alpha)-\sum_{j\in\mathcal J_n}(-1)^j\Phi(x_j,u_j)\right|
 \le\sum_{j\in I_n\mathbin{\triangle}\mathcal J_n}|\Phi(x_j,u_j)|,
\]
so the same tail argument applies in both directions of the symmetric
difference.  The probability that one of these costs exceeds
$\varepsilon L$ is $O_C(H/L)=o(1)$.  By integration of
\eqref{eq:Phi-tail}, the expected sum of the costs truncated at
$\varepsilon L$ is $O_C(H\log L)$; Markov's inequality after division by
$L$ makes this $o(1)$.  Finally let $C$ grow and use \eqref{eq:tau}.
\end{proof}

\begin{proof}[Proof of Theorem~\ref{thm:main}]
By Proposition~\ref{prop:principal}, Corollary~\ref{cor:principal-cauchy},
and Proposition~\ref{prop:remainder}, there are deterministic $c_n$ such that
\begin{equation}\label{eq:centered-final}
 \frac1L\sum_{j\in\mathcal J_n}(-1)^j\Phi(x_j,u_j)-c_n
 \Longrightarrow \operatorname{Cauchy}\!\left(0,\frac1{2\pi}\right).
\end{equation}
Lemma~\ref{lem:stopping} and Slutsky's theorem replace the bulk sum by
$S_n/L$.

It remains to remove the deterministic centering.  For almost every
$\alpha$,
\[
 \{k(1-\alpha)\}=1-\{k\alpha\},
 \qquad S_n(1-\alpha)=-S_n(\alpha).
\]
Thus $X_n=S_n/\log n$ has an exactly symmetric law.  The family
$X_n-c_n$ is tight by \eqref{eq:centered-final}; symmetry makes
$X_n+c_n$ tight as well.  Since their deterministic center difference is
$2c_n$, the sequence $(c_n)$ is bounded.  If $c_{n_k}\to c$, then
$X_{n_k}\Rightarrow Z+c$, where
$Z\sim\operatorname{Cauchy}(0,1/(2\pi))$.  The limit must be symmetric about
zero, and a nondegenerate translated Cauchy law has this symmetry only when
$c=0$.  Hence every convergent subsequence of $(c_n)$ tends to zero, so
$c_n\to0$.

The Cauchy distribution is continuous; therefore convergence in distribution
is equivalent to the asserted convergence of distribution functions.
\end{proof}

\section*{Acknowledgments}

The author thanks Ibby Mian and Shayaan Siddique of Millennium Research for
their detailed mathematical comments in an independent pre-collaboration audit
of an earlier version, published on 1 August 2026 \cite{MianSiddique26}, and
during the subsequent companion formalization.  Their comments led to
corrections and clarifications in Sections~3, 4, and~6.  Only after publication
of the audit did Mian, Siddique, and the author agree to collaborate separately
on a Lean~4 formalization of this proof.  Feedback arising from that distinct
companion project is collaborator feedback, not an additional independent
referee opinion; Mian and Siddique are not authors of this manuscript.

\section*{AI-use disclosure}

The author selected the problem and directed the iterative research and
prompting process.  The initial proposed answer, the proof architecture, and
the substantive mathematical steps were generated by OpenAI GPT-5.6 Sol Pro
and developed through iterative, human-directed AI collaboration.  OpenAI
Codex, under the author's direction, was used to assemble, edit, compile, and
perform automated consistency checks on this manuscript.  The author takes
responsibility for the submitted manuscript.

\begin{thebibliography}{99}

\bibitem{Baladi00}
V.~Baladi,
\emph{Positive Transfer Operators and Decay of Correlations},
World Scientific, Singapore, 2000.

\bibitem{Bloom1002}
T.~F. Bloom,
\emph{Erd\H{o}s Problem \#1002},
\url{https://www.erdosproblems.com/1002}, accessed July 13, 2026.

\bibitem{Borda25}
B.~Borda,
\emph{Equidistribution of continued fraction convergents in
$\operatorname{SL}(2,\mathbb Z_m)$ with an application to local discrepancy},
J. Mod. Dyn. 21 (2025), 327--359.

\bibitem{DS20}
D.~Dolgopyat and O.~Sarig,
\emph{Quenched and annealed temporal limit theorems for circle rotations},
Ast\'erisque 415 (2020), 59--85.

\bibitem{Er64b}
P. Erd\H{o}s,
\emph{Problems and results on diophantine approximations},
Compositio Math. 16 (1964), 52--65.

\bibitem{Kallenberg17}
O.~Kallenberg,
\emph{Random Measures, Theory and Applications},
Probability Theory and Stochastic Modelling 77,
Springer, Cham, 2017.

\bibitem{Kato95}
T.~Kato,
\emph{Perturbation Theory for Linear Operators},
Classics in Mathematics, Springer, Berlin, 1995.

\bibitem{Ke60}
H. Kesten,
\emph{Uniform distribution mod $1$},
Ann. of Math. (2) 71 (1960), 445--471.

\bibitem{MianSiddique26}
I.~Mian and S.~Siddique,
\emph{Independent verification of the Erd\H{o}s Problem 1002 claims},
Millennium Research report, 2026,
\url{https://github.com/ibrahimmian36/Tesserarius/blob/main/reports/erdos-1002.md}.

\bibitem{Vardi93}
I.~Vardi,
\emph{Dedekind sums have a limiting distribution},
Int. Math. Res. Not. 1993 (1993), no.~1, 1--12.

\bibitem{Zygmund59}
A.~Zygmund,
\emph{Trigonometric Series}, 2nd ed., Vols.~I--II,
Cambridge University Press, Cambridge, 1959.

\end{thebibliography}

\end{document}
